\section{Related Work}
\label{sec:related}

\subsection{Personalization of Text-to-Image Diffusion Models}
Text-to-image (T2I) diffusion models~\cite{ddpm,sd,wuerstchen} have emerged as state-of-the-art generative engines, pushing the boundaries of visual fidelity, compositionality, and prompt following. Their ability to synthesize high-dimensional data by learning the reverse drift of forward stochastic differential equations~\cite{songscore} has enabled broad creative applications. Beyond unconditional generation, research in diffusion personalization has expanded across multiple axes: text-guided localized image editing~\cite{prompt2prompt,instructpix2pix}, solving linear and non-linear inverse problems~\cite{dps}, concept-driven generation and subject identity preservation~\cite{dreambooth,lora}, and stylized synthesis~\cite{stylealigned,instantstyle,rout2025rbmodulation}. 

In exemplar-based personalization, the goal is to steer generation using visual prompts provided by reference images rather than exclusively relying on textual descriptions. To tailor a pre-trained diffusion model to a target visual style or subject, existing literature is bifurcated into two primary paradigms: (1) optimization-based fine-tuning (full-model or parameter-efficient) and (2) training-free test-time modulation. While fine-tuning adapts the underlying network weights, training-free methods dynamically alter the internal feature activations or sampling trajectories of frozen models.

\subsection{Fine-Tuning Diffusion Models for Style Personalization}
Fine-tuning approaches optimize network parameters to align internal representations with reference images. Full fine-tuning methods like DreamBooth~\cite{dreambooth} optimize the entire U-Net alongside unique text tokens, which offers high expressive capacity but demands significant storage and memory. To improve parameter efficiency, Low-Rank Adaptation (LoRA)~\cite{lora} injects trainable rank-decomposition matrices into cross-attention projections, vastly reducing checkpoint footprint. Building upon LoRA, ZipLoRA~\cite{ziplora} merges independently trained content and style LoRA weights via column-wise optimization, while StyleDrop~\cite{styledrop} employs iterative human-in-the-loop reward feedback to capture intricate visual styles. More recently, B-LoRA~\cite{b-lora} identifies distinct blocks within the diffusion architecture responsible for style versus content, training separate low-rank adapters. 

Despite their success, fine-tuning approaches exhibit severe practical limitations: (i) they depend heavily on multi-image collections (typically 4--10 images of the same consistent style), which are unavailable for single historical paintings or bespoke digital designs; (ii) executing backpropagation loops for every newly encountered style requires minutes of GPU compute, making them unsuitable for interactive workflows; and (iii) when forced to learn from only a \textit{single reference exemplar}, fine-tuning collapses into semantic memorization, overfitting to the depicted objects rather than extracting abstract artistic medium attributes.

\subsection{Training-Free Style Transfer and Attention Modulation}
The necessity to bypass training has catalyzed growing interest in training-free style personalization. Drawing inspiration from classical neural style transfer~\cite{Gatys_2016_CVPR,AdaIn}, early diffusion-based methods attempted to modulate intermediate activations. StyleAligned~\cite{stylealigned} synchronizes styles across a batch of generated images by sharing self-attention keys and values, normalizing queries and keys via Adaptive Instance Normalization (AdaIN). However, StyleAligned relies on deterministic DDIM inversion~\cite{DDIM} to cache reference keys and values, which introduces severe quantization drift and suppresses fine textural details; furthermore, it mandates coupled dual reverse diffusion processes that double GPU memory requirements. 

To eliminate inversion overhead, InstantStyle~\cite{instantstyle} injects visual features from a reference style image into heuristic subsets of cross-attention layers of a pre-trained IP-Adapter~\cite{ipadapter}. However, identifying effective injection layers is empirical and fails to generalize across distinct model backbones. Crucially, linear superposition of style features into attention layers directly induces catastrophic \textbf{content leakage}---inadvertently synthesizing foreign semantic objects from the style reference. To maintain layout fidelity, InstantStyle must condition on an external ControlNet~\cite{controlNet}, which suppresses prompt diversity and locks the output into rigid geometries. StyleID~\cite{styleid}, DiffStyler~\cite{diffstyler}, and FreeStyle~\cite{freestyle} manipulate self-attention layers to suppress leakage, yet frequently dilute stylistic expression to protect target geometry. 

Most recently, RB-Modulation~\cite{rout2025rbmodulation} introduced a training-free framework that bypasses external adapters and DDIM inversion. However, to suppress content leakage, RB-Modulation applies aggressive spatial mean-pooling over all style tokens ($\bar{\mathbf{t}}_{\text{style}} = \frac{1}{L} \sum \mathbf{t}_{\text{style}}^{(j)}$). While this collapses object identity, it eradicates spatial token variance, causing severe \textbf{under-stylization}, flat brushstrokes, and textural dilution. In contrast, \method overcomes this acute trade-off: through Blended Style Tokens with Null-space Orthogonal Projection, \method mathematically filters out semantic content interference without sacrificing the spatial variance vital for rich artistic texture.

\subsection{Guidance Mechanisms, Manifold Stability, and Stochastic Control}
Guidance techniques direct reverse diffusion trajectories toward desired data distributions. Classifier guidance~\cite{dhariwal2021diffusion} steers sampling via gradients of auxiliary classifiers, while Classifier-Free Guidance (CFG) blends conditional and unconditional score estimates. Tweedie's formula~\cite{tweedie} enables computing conditional expectations of clean data $\bx_0$ at intermediate timesteps, providing an analytical anchor for downstream guidance losses. 

RB-Modulation~\cite{rout2025rbmodulation} formulated style transfer as a stochastic optimal control problem using the Hamilton-Jacobi-Bellman (HJB) equation with a terminal style cost. However, computing raw gradient updates with respect to intermediate latents reveals a critical vulnerability: unconstrained guidance vectors possess substantial components antiparallel or orthogonal to the score function, driving diffusion trajectories off the valid data manifold and causing severe oversaturation artifacts and geometric collapse. Furthermore, RB-Modulation's reliance on the supervised Consistent Style Descriptor (CSD) leads to domain collapse when confronted with out-of-distribution artistic media. 

Moreover, foundational insights from SDEdit~\cite{sdedit} establish that the noise level in stochastic differential equations dictates a strict trade-off between input faithfulness and realism: early diffusion timesteps ($t \approx T$) govern coarse spatial layout and global semantics, whereas late timesteps resolve high-frequency textures. Uncontrolled late-stage guidance inevitably distorts geometry. \method synthesizes these insights: by coupling early-step clean latent pushforward (anchoring global color dynamics during layout formation) with Score-Orthogonal DINO guidance (constraining updates strictly to the score null space), \method achieves authentic brushstroke harmonization while rigorously preserving manifold stability.
